\documentclass[12pt,a4paper]{article}

\usepackage[T2A]{fontenc}
\usepackage[utf8]{inputenc}
\usepackage[russian]{babel}
\usepackage{geometry}
\usepackage{graphicx}
\usepackage{float}
\usepackage{hyperref}
\usepackage{booktabs}
\usepackage{enumitem}
\usepackage{xcolor}

\geometry{margin=2cm}
\hypersetup{
    colorlinks=true,
    linkcolor=blue,
    urlcolor=blue
}

\newcommand{\placeholderfigure}[2]{
    \begin{figure}[H]
        \centering
        \fbox{
            \parbox[c][0.28\textheight][c]{0.9\textwidth}{
                \centering
                \textbf{#1}\\[0.5em]
                Вставить скриншот: #2
            }
        }
        \caption{#1}
    \end{figure}
}

\title{Отчёт по мониторингу и профилированию веб-приложения}
\author{Студент: \underline{\hspace{6cm}}}
\date{\today}

\begin{document}
\maketitle

\section{Цель работы}
Целью работы является мониторинг Java EE-приложения при помощи VisualVM и профилировщика IDE, анализ пользовательских MBean, исследование использования памяти JVM, а также локализация и устранение проблемы производительности.

\section{Исходные данные}
В рамках лабораторной работы в приложении уже были реализованы два пользовательских MBean:
\begin{itemize}[leftmargin=*]
    \item \texttt{PointStatistics} --- считает общее число установленных точек, число промахов и число промахов подряд;
    \item \texttt{ClickInterval} --- считает средний интервал между кликами пользователя по координатной плоскости.
\end{itemize}

Для мониторинга использовались утилиты \texttt{VisualVM} и встроенный профилировщик IDE. Приложение запускалось на сервере \texttt{<указать сервер и версию>} с JVM \texttt{<указать JVM и версию>}.

\section{Мониторинг в VisualVM}
\subsection{Подключение к приложению}
Для подключения к приложению в VisualVM был выполнен следующий порядок действий:
\begin{enumerate}[leftmargin=*]
    \item Запустить сервер приложений с развернутым \texttt{WAR}-архивом.
    \item Открыть \texttt{VisualVM} и дождаться появления процесса сервера в разделе \texttt{Local}.
    \item Выбрать процесс JVM, в котором выполняется приложение.
    \item Перейти на вкладку \texttt{MBeans} и убедиться, что зарегистрированы пользовательские объекты \texttt{programming.itmo:type=PointStatistics} и \texttt{programming.itmo:type=ClickInterval}.
\end{enumerate}

\subsection{Наблюдение за MBean}
Для получения изменяющихся показаний был выполнен сценарий:
\begin{enumerate}[leftmargin=*]
    \item Несколько раз отправить точки через форму.
    \item Несколько раз кликнуть по координатной плоскости.
    \item Сформировать последовательность из трёх промахов подряд.
    \item Наблюдать изменение атрибутов MBean во времени.
\end{enumerate}

В ходе наблюдения должны изменяться следующие показатели:
\begin{itemize}[leftmargin=*]
    \item \texttt{PointStatistics.TotalPoints};
    \item \texttt{PointStatistics.MissedPoints};
    \item \texttt{PointStatistics.ConsecutiveMisses};
    \item \texttt{ClickInterval.ClickCount};
    \item \texttt{ClickInterval.IntervalCount};
    \item \texttt{ClickInterval.AverageIntervalMillis}.
\end{itemize}

\placeholderfigure{График изменения показаний MBean}{графики из вкладки MBeans / Monitor в VisualVM}

\subsection{Анализ памяти JVM}
Для анализа памяти использовалась вкладка \texttt{Sampler} или \texttt{Profiler} в режиме памяти:
\begin{enumerate}[leftmargin=*]
    \item Открыть анализ памяти для процесса приложения.
    \item Выполнить пользовательский сценарий с добавлением большого числа точек.
    \item Снять heap histogram или список классов по объёму занятой памяти.
    \item Найти класс, экземпляры которого суммарно занимают наибольший объём памяти.
    \item Установить, в экземплярах какого пользовательского класса содержатся объекты этого типа.
\end{enumerate}

\noindent
\textbf{Результат анализа памяти:}
\begin{itemize}[leftmargin=*]
    \item класс, занявший наибольший объём памяти: \texttt{<заполнить по VisualVM>};
    \item пользовательский класс, в экземплярах которого находятся эти объекты: \texttt{<заполнить по VisualVM>}.
\end{itemize}

\placeholderfigure{Результаты анализа памяти JVM}{скриншот heap histogram / memory sampler}

\section{Локализация проблемы производительности}
\subsection{Описание выявленной проблемы}
В исходной версии приложения производительная проблема проявлялась при изменении параметра \texttt{R}:
\begin{itemize}[leftmargin=*]
    \item при каждом изменении поля \texttt{R} выполнялся AJAX-запрос на сервер;
    \item серверный метод \texttt{PointsBean.updateFilteredPoints()} каждый раз заново считывал все точки из базы данных;
    \item после чтения всех точек выполнялись повторный пересчёт попадания для каждой точки и полная сериализация списка в JSON;
    \item при наборе значения \texttt{R} это приводило к серии частых одинаковых запросов и лишней нагрузке на CPU и БД.
\end{itemize}

Следствием были лишние вызовы методов доступа к данным, рост числа коротких серверных запросов и избыточные затраты времени на обновление графика.

\subsection{Как проблема была локализована}
Для локализации использовался следующий алгоритм:
\begin{enumerate}[leftmargin=*]
    \item Запустить приложение и подключиться к его JVM из \texttt{VisualVM}.
    \item Открыть профилирование CPU.
    \item Начать запись профиля.
    \item Многократно изменять значение \texttt{R} в форме, не отправляя новые точки.
    \item Остановить профилирование и отсортировать методы по Self Time / Total Time.
    \item Проверить, что среди горячих методов находятся:
    \begin{itemize}[leftmargin=*]
        \item \texttt{programming.itmo.beans.PointsBean.updateFilteredPoints};
        \item \texttt{programming.itmo.repositories.PointRepositoryORM.getAllPoints};
        \item методы сериализации JSON и создания объектов \texttt{PointDTO}.
    \end{itemize}
    \item Повторить тот же сценарий во встроенном профилировщике IDE для подтверждения результатов.
\end{enumerate}

\placeholderfigure{Локализация горячих методов в VisualVM}{скриншот CPU profiler с методом updateFilteredPoints}

\placeholderfigure{Подтверждение в профилировщике IDE}{скриншот hotspot / call tree из IntelliJ IDEA, NetBeans или Eclipse}

\subsection{Внесённые изменения}
Для устранения проблемы были выполнены следующие изменения:
\begin{itemize}[leftmargin=*]
    \item в \texttt{PointsBean} добавлен кэш списка исторических точек в рамках сессии пользователя;
    \item метод обновления графика переведён на работу с кэшированными точками вместо повторного чтения всей таблицы БД при каждом изменении \texttt{R};
    \item для отрисовки графика создаётся отдельная проекция точек с пересчитанным признаком попадания, чтобы не портить исходные данные таблицы;
    \item в \texttt{draw.js} добавлена debounce-синхронизация с сервером, чтобы не отправлять каскад AJAX-запросов при быстром наборе \texttt{R};
    \item серверная валидация поля \texttt{R} переведена с события \texttt{keyup} на \texttt{blur}, так как интерактивная перерисовка графика уже выполняется на клиенте.
\end{itemize}

\subsection{Файлы, в которых выполнены изменения}
\begin{itemize}[leftmargin=*]
    \item \texttt{src/main/java/programming/itmo/beans/PointsBean.java};
    \item \texttt{src/main/java/programming/itmo/util/CheckAreaUtil.java};
    \item \texttt{src/main/webapp/main.xhtml};
    \item \texttt{src/main/webapp/resources/scripts/draw.js}.
\end{itemize}

\subsection{Описание путей устранения проблемы}
Основной путь устранения заключался не в оптимизации одной операции, а в сокращении числа бессмысленных повторений:
\begin{itemize}[leftmargin=*]
    \item исключено повторное чтение одной и той же коллекции точек из БД при каждом изменении \texttt{R};
    \item уменьшено количество серверных обращений при вводе значения \texttt{R};
    \item сохранено текущее пользовательское поведение интерфейса: область продолжает перерисовываться сразу, а данные для графика приходят с сервера заметно реже.
\end{itemize}

\subsection{Проверка результата}
После внесения изменений необходимо повторить исходный сценарий профилирования:
\begin{enumerate}[leftmargin=*]
    \item снова запустить CPU profiling;
    \item повторить многократное изменение \texttt{R};
    \item сравнить число вызовов и суммарное время работы методов до и после исправления;
    \item зафиксировать уменьшение нагрузки на методы доступа к данным и обновления списка точек.
\end{enumerate}

\noindent
\textbf{Итог сравнения до и после исправления:}
\begin{itemize}[leftmargin=*]
    \item количество вызовов \texttt{PointRepositoryORM.getAllPoints}: \texttt{<до>} $\rightarrow$ \texttt{<после>};
    \item суммарное время в \texttt{PointsBean.updateFilteredPoints}: \texttt{<до>} $\rightarrow$ \texttt{<после>};
    \item субъективная реакция интерфейса: \texttt{<описать после замеров>}.
\end{itemize}

\placeholderfigure{Сравнение профиля до и после исправления}{скриншоты hotspot до и после оптимизации}

\section{Воспроизводимость}
Для демонстрации преподавателю процесс можно воспроизвести следующим образом:
\begin{enumerate}[leftmargin=*]
    \item запустить сервер приложений с развернутым проектом;
    \item открыть \texttt{VisualVM} и подключиться к JVM сервера;
    \item показать вкладку \texttt{MBeans} и изменение пользовательских атрибутов во времени;
    \item открыть анализ памяти и показать класс, занимающий наибольший объём;
    \item запустить CPU profiling;
    \item воспроизвести сценарий многократного изменения \texttt{R};
    \item показать, какие методы были горячими до исправления и как изменился профиль после исправления.
\end{enumerate}

\section{Вывод}
В ходе работы были исследованы пользовательские MBean приложения, выполнен анализ использования памяти JVM, локализована проблема производительности, связанная с повторными запросами и лишней обработкой данных при изменении параметра \texttt{R}, и предложено исправление, уменьшающее нагрузку на серверную часть приложения.

\end{document}
